\chapter{Measuring and Time}\label{ch:g1:measure}

Which pencil is longer? What day is it? How much does a sweet cost?
This chapter compares lengths and measures them in centimeters,
puts the days of the week in order, reads whole hours on the clock,
and counts small coins.

\section{Longer and shorter}

\begin{method}[Comparing lengths]\label{met:g1:measure:compare}
To compare two lengths, line them up with \emph{one end matched}:
\begin{enumerate}
  \item put both objects side by side, starting from the same line;
  \item the one that reaches further is \emph{longer}, the other is
        \emph{shorter};
  \item if both reach the same place, they are the \emph{same
        length}.
\end{enumerate}
Matching the start is the whole secret --- comparing from crooked
starts fools the eye.
\end{method}

\begin{omfigure}
\begin{tikzpicture}[scale=0.9]
  \draw[very thick, omDef] (0,0.6) -- (3.4,0.6);
  \draw[very thick, omThm] (0,0) -- (2.2,0);
  \draw[gray, dashed] (0,-0.2) -- (0,0.8);
  \node[right, font=\small] at (3.6,0.6) {longer};
  \node[right, font=\small] at (3.6,0) {shorter};
\end{tikzpicture}

{\small Two sticks, matched at the left: the blue one reaches
further, so it is longer.}
\end{omfigure}

\begin{definition}[The centimeter]\label{def:g1:measure:cm}
To say \emph{how} long, we measure with a ruler in
\emph{centimeters} (cm). Put the $0$ of the ruler at the start of
the object, and read the number at its end.
\end{definition}

\begin{omfigure}
\begin{tikzpicture}[scale=1.0]
  \draw[thick] (0,0) rectangle (8,0.7);
  \foreach \i in {0,...,8} {
    \draw[thick] (\i,0.7) -- (\i,0.4);
    \node[font=\footnotesize] at (\i,0.2) {\i};
  }
  \draw[very thick, omThm] (0,1.0) -- (5,1.0);
  \node[omThm, above, font=\small] at (2.5,1.0) {the pencil: $5$ cm};
\end{tikzpicture}

{\small Measuring: the $0$ at the start, read the end. This pencil
is $5$ cm long. (Millimeters and meters come in
\cref{ch:g3:measure}.)}
\end{omfigure}

\section{The days of the week}

\begin{example}[The seven days]\label{ex:g1:measure:days}
The week has $7$ days, always in the same order:
\[
\text{Monday, Tuesday, Wednesday, Thursday, Friday, Saturday,
Sunday.}
\]
After Sunday, Monday starts again. \emph{Yesterday} is the day
before, \emph{tomorrow} the day after: if today is Wednesday,
yesterday was Tuesday and tomorrow is Thursday.
\end{example}

\section{Whole hours on the clock}

\begin{method}[Reading a whole hour]\label{met:g1:measure:clock}
When the long hand points straight up to the $12$, it is a whole
hour: read the number the \emph{short} hand points to. Short hand on
the $3$, long hand on the $12$: it is $3$ o'clock.
\end{method}

\begin{omfigure}
\begin{tikzpicture}[scale=0.95]
  \draw[thick] (0,0) circle (1.3);
  \foreach \h in {1,...,12} {
    \node[font=\footnotesize] at ({90-\h*30}:1.05) {\h};
    \draw ({90-\h*30}:1.22) -- ({90-\h*30}:1.3);
  }
  \draw[very thick, omDef] (0,0) -- (0:0.62);
  \draw[thick, omThm] (0,0) -- (90:1.02);
  \node[below, font=\small] at (0,-1.55) {$3$ o'clock};
  \begin{scope}[xshift=4.4cm]
  \draw[thick] (0,0) circle (1.3);
  \foreach \h in {1,...,12} {
    \node[font=\footnotesize] at ({90-\h*30}:1.05) {\h};
    \draw ({90-\h*30}:1.22) -- ({90-\h*30}:1.3);
  }
  \draw[very thick, omDef] (0,0) -- ({90-8*30}:0.62);
  \draw[thick, omThm] (0,0) -- (90:1.02);
  \node[below, font=\small] at (0,-1.55) {$8$ o'clock};
  \end{scope}
\end{tikzpicture}

{\small Whole hours: the long red hand on the $12$, the short blue
hand giving the hour. (Half and quarter hours come in
\cref{ch:g2:measure}.)}
\end{omfigure}

\section{Small coins}

\begin{example}[Making an amount]\label{ex:g1:measure:coins}
With coins of $1$, $2$ and $5$, many amounts can be made in several
ways. To make $7$:
\[
5 + 2, \qquad
5 + 1 + 1, \qquad
2 + 2 + 2 + 1 .
\]
The biggest coins first is usually the quickest way, using the
fewest coins.
\end{example}

\begin{omfigure}
\begin{tikzpicture}[scale=0.85]
  \draw[thick, fill=omMeth!20] (0,0) circle (0.5);
  \node[font=\small] at (0,0) {$5$};
  \node[font=\large] at (1.0,0) {$+$};
  \draw[thick, fill=omDef!20] (2.0,0) circle (0.42);
  \node[font=\small] at (2.0,0) {$2$};
  \node[font=\large] at (3.0,0) {$=$};
  \node[font=\small] at (4.0,0) {$7$};
\end{tikzpicture}

{\small A five-coin and a two-coin make seven.}
\end{omfigure}

\section{Exercises}

\begin{exercise}[$\star$]\label{exo:g1:measure:1}
Line up three pencils matched at one end. Which is the longest?
Which is the shortest?
\end{exercise}

\begin{exercise}[$\star$]\label{exo:g1:measure:2}
Measure with your ruler: the long side of this book; your little
finger; a spoon. Answer in whole centimeters.
\end{exercise}

\begin{exercise}[$\star$]\label{exo:g1:measure:3}
Name the day: after Monday; before Sunday; two days after Friday.
\end{exercise}

\begin{exercise}[$\star$]\label{exo:g1:measure:4}
Today is Thursday. What day was yesterday? What day is tomorrow?
What day is in three days?
\end{exercise}

\begin{exercise}[$\star$]\label{exo:g1:measure:5}
Draw a clock showing $6$ o'clock. Where is the short hand? Where is
the long hand?
\end{exercise}

\begin{exercise}[$\star$]\label{exo:g1:measure:6}
Read the whole hour: short hand on $5$, long hand on $12$; short
hand on $11$, long hand on $12$.
\end{exercise}

\begin{exercise}[$\star$]\label{exo:g1:measure:7}
Make $8$ with coins of $1$, $2$ and $5$, in two different ways.
\end{exercise}

\begin{exercise}[$\star$]\label{exo:g1:measure:8}
Leo has a $5$ coin and three $2$ coins. How much money has he
altogether?
\end{exercise}

\begin{exercise}[$\star\star$]\label{exo:g1:measure:9}
A sweet costs $3$. Zoe pays with a $5$ coin. How much change does
she get back? (Count up from $3$ to $5$.)
\end{exercise}
